\documentclass[12pt, a4paper]{article}

% ── 中文支持（Overleaf 请选 XeLaTeX 编译器）──────────────────────────────────
\usepackage[UTF8, heading=true]{ctex}

% ── 页面布局 ─────────────────────────────────────────────────────────────────
\usepackage[top=2.5cm, bottom=2.5cm, left=3cm, right=3cm]{geometry}

% ── 数学 ─────────────────────────────────────────────────────────────────────
\usepackage{amsmath}
\usepackage{amssymb}
\usepackage{bm}

% ── 图表 ─────────────────────────────────────────────────────────────────────
\usepackage{graphicx}
\usepackage{float}
\usepackage{booktabs}
\usepackage{multirow}
\usepackage{tabularx}
\usepackage{array}
\usepackage{caption}
\usepackage{subcaption}

% ── 代码块 ───────────────────────────────────────────────────────────────────
\usepackage{listings}
\usepackage{xcolor}
\lstset{
    basicstyle=\small\ttfamily,
    breaklines=true,
    frame=single,
    rulecolor=\color{gray},
    commentstyle=\color{gray},
    keywordstyle=\color{blue},
    stringstyle=\color{orange},
    numbers=left,
    numberstyle=\tiny\color{gray},
    tabsize=2,
    showstringspaces=false,
}

% ── 超链接 ───────────────────────────────────────────────────────────────────
\usepackage[colorlinks=true, linkcolor=blue, citecolor=blue, urlcolor=blue]{hyperref}

% ── 行距 ─────────────────────────────────────────────────────────────────────
\usepackage{setspace}
\onehalfspacing

% ── 参考文献 ─────────────────────────────────────────────────────────────────
\usepackage[numbers, sort&compress]{natbib}

% ── 标题信息 ─────────────────────────────────────────────────────────────────
\title{\textbf{面向航空发动机装配的多模态混合检索增强生成框架}}
\author{作者姓名}
\date{\today}

% ═════════════════════════════════════════════════════════════════════════════
\begin{document}

\maketitle

% ── 摘要 ─────────────────────────────────────────────────────────────────────
\begin{abstract}
针对航空发动机装配领域知识密集、文档多模态、查询复杂的特点，
本文提出一种多模态混合检索增强生成（Multimodal Hybrid RAG）框架。
% TODO: 上传 chapters/abstract.tex 后替换为 \input{chapters/abstract}
\end{abstract}

\newpage
\tableofcontents
\newpage

% ═══════════════════════════════════════════════════════════════════════════
% 第 1 章  引言
% ═══════════════════════════════════════════════════════════════════════════
\section{引言}

\subsection{研究背景与动机}
% TODO

\subsection{现有方法局限性}
% TODO

\subsection{本文主要贡献}
% TODO

\subsection{论文结构}
% TODO

% 完成后取消注释：% ═════════════════════════════════════════════════════════════════════════════
%  第一章  引言
% ═════════════════════════════════════════════════════════════════════════════
\section{引言}
\label{sec:intro}

% ─────────────────────────────────────────────────────────────────────────────
\subsection{研究背景与动机}
\label{sec:background}

航空发动机是现代航空工业的核心动力装置，其装配过程涉及数千个零部件、
数百道工序以及严格的公差配合约束，是典型的知识密集型、高精度制造场景。
支撑装配作业的知识分散在三类异构数据源中：
\textbf{维修手册}以 PDF 文档形式承载工序逻辑、工具要求和技术规范等过程性知识；
\textbf{BOM（物料清单）表}以结构化电子表格形式提供零件编号、名称和数量等身份信息；
\textbf{三维CAD模型}以几何数据形式编码零件装配层级关系、配合约束和几何公差等空间结构知识。
三类数据源各有其信息边界：手册不含装配结构信息，BOM 不含装配步骤，CAD 不含操作文字描述——
单一数据源均无法构成完整的装配知识底座。
知识图谱（Knowledge Graph）作为结构化知识表示的核心工具，
以三元组形式编码实体间的语义关系，已被证明在工业领域知识管理中具有重要价值
\cite{hogan2021kg,ji2022kgsurvey}。
如何融合三源异构知识，并以自然语言交互方式辅助装配工人完成精密操作，
是提升航空制造智能化水平的重要课题。

近年来，大语言模型（Large Language Model，LLM）在自然语言理解与生成任务上取得了重大突破，
但其固有的三类局限制约了在专业领域的可靠应用：
\textbf{参数知识幻觉}（生成听似合理的错误事实）、
\textbf{训练知识截止}（无法获取最新领域数据），
以及\textbf{专业推理能力不足}（在涉及精确数值和多跳逻辑的专业领域问题上表现不稳定）
\cite{zhao2023llmsurvey}。
检索增强生成（Retrieval-Augmented Generation，RAG）技术为上述问题提供了系统性解决路径
\cite{lewis2020rag}：
通过在推理时动态检索外部知识库，
将相关文本片段注入LLM上下文窗口，
使模型的生成过程以外部事实为锚，
而非单纯依赖参数记忆。
RAG技术自提出以来经历了从Naive RAG到Advanced RAG再到GraphRAG的快速演进
\cite{gao2023ragsurvey}，
在医疗、法律、金融等专业领域已取得显著进展。

然而，现有RAG方法多针对单一文本知识源设计，
难以满足航空发动机装配领域的以下三类特殊需求：

\begin{enumerate}
  \item \textbf{精确标识符匹配}：
        零件编号（如 PN: 1002-AX）、公差数值（如"0.05\textasciitilde0.12\,mm"）、
        扭矩规格等精确标识符对检索准确性至关重要，
        而单纯的密集向量检索因词汇平滑而弱化此类精确匹配能力；
  \item \textbf{结构化推理需求}：
        零件组成关系、装配工序依赖、配合约束等知识天然具有图结构，
        需要多跳关系推理，远超向量相似度检索的能力边界；
  \item \textbf{跨源一致性约束}：
        手册、BOM 与 CAD 对同一零件或工序常采用不同命名方式，
        若缺乏统一实体注册与跨源对齐，系统容易在装配层级、工具要求和技术规范之间产生断裂；
        此外，检索噪声（上下文中混入的不相关段落）会导致LLM生成质量显著下降
        \cite{liu2023lostinmiddle}，
        对高精度装配场景危害尤甚。
\end{enumerate}

现有工作已开始关注RAG方法在工业故障诊断与装配引导场景中的应用。
在航空故障诊断领域，Tang等人\cite{tang2023aircraft}和Meng等人\cite{meng2024aircraft}
分别系统综述并实现了面向航空故障诊断的知识图谱平台；
Wu等人\cite{wu2024engine}构建了专门面向航空发动机润滑系统的故障KG，
证明了KG在航空发动机专业场景下的工程价值。
Liu等人\cite{liu2024jointkg}在航空装配功能测试场景下，
将知识图谱通过前缀微调嵌入大语言模型，实现了98.5\%的故障诊断准确率，
证明了专业领域KG与LLM协同推理的有效性。
Liu等人\cite{liu2023avkg}构建了含5类实体、3类关系共1308组三元组的
航空装配知识图谱，并与Neo4j图数据库集成，为装配知识管理奠定了基础。
Shi等人\cite{shi2022arkg}提出基于"功能-结构-工序-装配资源"多维本体的
装配资源知识图谱（ARKG），将CAD模型数据经本体映射后融入知识图谱，
支持工程决策的知识共享与复用。
Shen等人\cite{shen2025hybridrag}提出了HybridRAG框架，
将知识图谱（KG）与向量检索融合，
并在飞机故障排查场景中验证了混合检索相较于单一稠密检索路径的优越性，
F1提升约4\%、幻觉率降低约7\%。
然而，上述工作多以手册文本或 CAD 属性的单一侧面为主，
尚未将 BOM 表的权威零件身份信息、CAD/IPC 装配结构信息与维修手册工序知识系统融入同一图谱。
这使其在回答"某零件属于哪个上级装配体"、"某工序需要哪些工具和规范"等问题时，
容易出现实体缺失、关系链断裂和跨章节召回不稳定等现象。
现有方法在装配这一三源知识协同的精细场景下仍存在明显局限。

% ─────────────────────────────────────────────────────────────────────────────
\subsection{现有方法局限性}
\label{sec:limitation}

Fan等人\cite{fan2024ragsurvey}的综述指出，RAG方法的核心挑战在于检索质量、
增强策略和生成可靠性三个维度；Gao等人\cite{gao2023ragsurvey}进一步将RAG发展
划分为Naive RAG、Advanced RAG和GraphRAG三个演进阶段，指出GraphRAG在结构化推理
方面具有显著优势但知识来源仍局限于文本。
结合航空发动机装配的特殊需求，本文进一步识别出以下三方面不足：

\textbf{（1）缺乏稀疏-稠密互补机制。}
大多数方案选择单一的稠密向量检索路径，
在处理含精确零件编号和数字标识符的领域特定查询时召回率偏低。
BM25等稀疏检索方法虽对词汇精确匹配有天然优势，
但单独使用时语义泛化能力不足，
两类方法的互补融合尚未在装配领域得到充分研究。

\textbf{（2）三源异构知识融合不足。}
现有工业知识图谱工作多从维修手册文本单一来源提取三元组：
Liu等人\cite{liu2023avkg}的航空装配KG仅含5类实体和3类关系，
Shi等人\cite{shi2022arkg}的ARKG虽引入CAD模型数据，但未实现BOM、CAD、手册三者的
系统融合与实体对齐。三源数据的互补边界清晰
（手册提供过程知识，BOM提供身份信息，CAD提供几何拓扑），
但现有方法缺乏对三者进行统一本体建模和跨源实体对齐的完整框架。

\textbf{（3）多源异质结果融合方法简单。}
多路检索的结果评分处于不同尺度空间，
现有方法多采用线性加权融合（$\text{IR} = \alpha \cdot \text{score}_{\text{dense}} + \beta \cdot \text{score}_{\text{sparse}}$）
\cite{shen2025hybridrag}，
易受评分分布差异影响，
缺乏对文本块与图谱子图等异质内容统一排序的无参融合机制。

% ─────────────────────────────────────────────────────────────────────────────
\subsection{本文主要贡献}
\label{sec:contribution}

针对上述不足，本文提出面向航空发动机装配的
\textbf{多源混合检索增强生成}框架
（Multi-source Hybrid RAG，MHRAG），
主要贡献如下：

\begin{enumerate}
  \item \textbf{三源知识图谱构建方法}。
        设计了以BOM解析、CAD结构化解析、LLM手册三元组抽取为三轨的
        多智能体协同知识抽取流水线，三轨结果通过实体对齐后
        汇入统一的Neo4j图数据库。
        本体模式包含7类实体（含 \textit{GeoConstraint} 几何约束节点）
        和9类关系，清晰界定了三源数据的语义边界：
        BOM提供零件身份注册表，CAD提供装配层级与几何约束，
        手册提供工序过程与技术规范。
        知识融合阶段引入验证智能体进行一致性校验，
        有效解决同名异物和信息冲突问题。

  \item \textbf{三路并行混合检索架构}。
        在密集向量检索（BGE-M3）和BM25稀疏检索之外，
        引入基于BGE-M3语义近邻搜索（ANN）的知识图谱子图召回路径，
        使系统同时具备语义召回、精确字符串匹配和结构化关系补全能力。
        三路结果通过互惠排名融合（RRF）统一排序，
        无需对各路评分进行归一化处理，对单路检索失败具有天然鲁棒性；
        相比传统的关键词字符串匹配（CONTAINS）KG检索方法，
        语义ANN方案更适合处理中英混合查询与专业术语别名。

  \item \textbf{面向装配引导的生成与交互设计}。
        引入多查询扩展策略缓解词汇不匹配问题；
        采用基于BERT的交叉编码器精排进一步提升上下文质量；
        设计结构化系统提示使LLM以"逐步确认"模式输出操作指令，
        符合航空装配的安全性要求；
        答案通过SSE流式返回并附带来源引用，
        支持前端展示可追溯的手册片段与图谱关系。

  \item \textbf{系统实现与验证}。
        基于FastAPI + Neo4j + Qdrant实现了完整的后端服务，
        前端以React实现交互式问答界面，
        在 PT6A 整机维修手册（10 个 ATA 章节、约 1.78\,MB Markdown 文本）上构建出
        含 7{,}333 节点 / 8{,}752 边的知识图谱。
        在含 60 题精评（72-30 压气机子系统）与 20 题跨章节粗评的双层 QA 评测集上，
        本文方法（MHRAG）在 Tier-1 精评的 LLM-Judge Relevance、Completeness 与
        Correctness 维度取得 3.50、3.02 与 2.35 的表现；
        在 Tier-2 跨章节粗评中 Completeness 与 Correctness 维度领先所有单路基线
        （相对 BM25-only 分别提升 8.9\% 与 1.9\%）；
        消融实验进一步分析各路径的独立贡献：
        KG 路径主要提升答案完整性与事实一致性，
        BM25 路径则保证零件编号、扭矩和公差等精确表达的可召回性。
        工程性能维度，本文方法在 7.76\,s 端到端延迟和 1.60\,GB 本地显存占用下
        取得 35.00\% 的答非所问率，
        契合消费级 GPU 的工业部署约束。
\end{enumerate}

% ─────────────────────────────────────────────────────────────────────────────
\subsection{论文结构}
\label{sec:structure}

本文其余部分组织如下：
第\ref{sec:related}节回顾RAG、知识图谱构建与混合检索的相关工作；
第\ref{sec:method}节详细介绍本文提出的MHRAG框架各组件；
第\ref{sec:experiment}节描述实验配置、数据集构建与评价指标，
并汇报与基线方法的对比结果及消融分析；
第\ref{sec:discussion}节讨论各检索路径的贡献、知识图谱质量影响以及方法局限性；
第\ref{sec:conclusion}节总结全文并展望未来工作方向。



% ═══════════════════════════════════════════════════════════════════════════
% 第 2 章  相关工作
% ═══════════════════════════════════════════════════════════════════════════
\section{相关工作}

\subsection{检索增强生成（RAG）综述}
% TODO

\subsection{知识图谱构建与推理}
% TODO

\subsection{多模态信息检索}
% TODO

\subsection{航空发动机装配智能化}
% TODO

% 完成后取消注释：% ═════════════════════════════════════════════════════════════════════════════
%  第二章  相关工作
% ═════════════════════════════════════════════════════════════════════════════
\section{相关工作}
\label{sec:related}

% ─────────────────────────────────────────────────────────────────────────────
\subsection{检索增强生成（RAG）综述}
\label{sec:related-rag}

\subsubsection{RAG技术演进}

检索增强生成（Retrieval-Augmented Generation，RAG）由Lewis等人\cite{lewis2020rag}
提出，其核心思想是在LLM推理阶段动态检索外部知识库，
将召回的相关文本片段注入上下文窗口，引导模型生成更准确、更具事实依据的答案。
Fan等人\cite{fan2024ragsurvey}从架构设计、训练策略与应用场景三个维度
对RAG-LLMs研究进行了全面梳理，认为RAG在开放域问答等知识密集型任务中展现出显著潜力。
Gao等人\cite{gao2023ragsurvey}的综述将RAG的发展划分为三个递进阶段：

\begin{figure}[htbp]
  \centering
  \includegraphics[width=\linewidth]{figures/fig7_rag_evolution.pdf}
  \caption{RAG技术三阶段演进对比：Naive RAG、Advanced RAG与GraphRAG}
  \label{fig:rag-evolution}
\end{figure}

\textbf{Naive RAG}是最早期范式，以稠密向量检索为核心，
直接将Top-$K$文档块拼接入提示词。
该范式在通用问答任务上表现有效，
但在专业领域面临检索精度低、上下文噪声大等问题\cite{fan2024ragsurvey}。
Liu等人\cite{liu2023lostinmiddle}通过实验发现，
当无关段落混入检索上下文时，LLM的生成质量会显著下降——
这一"中间丢失"（Lost in the Middle）现象在长上下文场景尤为突出，
揭示了Naive RAG在精确查询场景中的根本局限。

\textbf{Advanced RAG}在上述基础上引入查询改写、混合检索与精排等优化模块，
以多层次干预系统性提升检索质量。
代表性改进包括：
多查询扩展\cite{ma2023query}（用LLM将原始查询改写为多种表述后并行检索）、
稀疏-稠密混合检索\cite{luan2021sparse}（BM25与向量检索互补融合）、
以及交叉编码器精排\cite{nogueira2019passage}（在粗排后进行细粒度语义重排）。

\textbf{GraphRAG}则走向了更深层的结构化方向。
Edge等人\cite{edge2024graphrag}提出将知识图谱与RAG结合，
以图结构表示实体关系，支持多跳推理，
有效弥补了向量检索在结构化推理方面的能力边界。
本文在GraphRAG框架的基础上，针对航空装配场景的特殊需求，
进一步设计了三源知识图谱与基于密集语义近邻搜索（ANN）的子图召回路径——
通过将BGE-M3统一编码用户查询与图谱实体描述，使中英混合查询能够直接命中
图谱中的语义近邻锚点，再以一跳邻域扩展提取相关子图，
解决了纯文本图谱无法覆盖几何约束知识、且关键词匹配对中文查询召回率偏低的双重问题。

\subsubsection{稀疏-稠密混合检索}

早期RAG工作以双编码器（Bi-Encoder）稠密检索为主，
代表性模型包括DPR\cite{karpukhin2020dpr}和BGE系列\cite{bge2023}。
稠密检索擅长捕捉语义相似性，但在面对零件编号（PN）等精确标识符时匹配能力不足。

BM25\cite{robertson2009bm25}作为经典概率检索模型，
在词汇层面具有天然的精确匹配优势。
研究表明，稀疏检索与稠密向量检索互补融合后，
综合召回率和准确率通常优于任一单路方法\cite{luan2021sparse}。
现有工业RAG工作（如HybridRAG\cite{shen2025hybridrag}）采用
线性加权融合（$\text{score} = \alpha \cdot s_{\text{dense}} + \beta \cdot s_{\text{sparse}}$），
该方式依赖人工调整权重且对评分分布敏感。
本文改用互惠排名融合（RRF）\cite{cormack2009rrf}，
无需对不同评分空间进行归一化处理，对单路检索失败具有天然鲁棒性，
且可自然地将文本块与图谱子图等异质内容统一排序。

精排阶段，交叉编码器（Cross-Encoder）基于Transformer编码器
\cite{vaswani2017attention,devlin2019bert}，
在推理时对查询与候选段落进行全注意力交互，
相比双编码器的独立编码范式能更精准地捕捉细粒度语义对齐关系\cite{nogueira2019passage}。

% ─────────────────────────────────────────────────────────────────────────────
\subsection{知识图谱构建与推理}
\label{sec:related-kg}

知识图谱以三元组~$(e_h, r, e_t)$~的形式表示实体间结构化关系，
为需要多跳推理的复杂问答提供向量检索难以胜任的能力。
Hogan等人\cite{hogan2021kg}对知识图谱的表示方法、构建技术与应用进行了权威综述，
指出KG在知识密集型任务中具有重要的结构化推理优势。
Ji等人\cite{ji2022kgsurvey}从表示学习、知识获取与应用三个维度系统梳理了KG研究进展，
为本文的本体设计和实体对齐方法奠定了理论基础。
Pan等人\cite{pan2024kgllmsurvey}进一步梳理了知识图谱与LLM融合的路线图，
认为两者结合在专业领域问答中具有显著协同优势——
这也是本文选择KG+RAG融合路线的核心依据之一。

\subsubsection{基于LLM的零样本三元组抽取}

早期命名实体识别（NER）与关系抽取（RE）工作依赖大规模标注数据，
在航空制造这类低资源领域面临标注成本高昂的瓶颈。
近年来，Pan 等\cite{pan2024kgllmsurvey}的综述系统梳理了
基于预训练语言模型的零样本或少样本三元组抽取范式——
通过将本体模式编码进少样本提示词（Few-shot Prompting）驱动 LLM 直接输出结构化三元组——
可大幅降低知识图谱的构建成本，使低资源领域的 KG 构建成为可能。

为进一步提升抽取质量，Edge等人\cite{edge2024graphrag}提出的Gleaning机制
在首轮抽取后以第二轮提示要求模型补充遗漏实体；
实体对齐（Entity Linking）则通过缩写词典映射、编辑距离模糊匹配和语义向量
三级级联方式\cite{wu2020corefresolution}解决同义表达冲突，
从而保证图谱内部的一致性。

\subsubsection{基于CAD几何数据的结构化知识抽取}

现有工业知识图谱研究多聚焦于文本、数据库和传感器数据，
鲜有工作直接从三维CAD模型中抽取精确的结构化知识。
Shi等人\cite{shi2022arkg}提出基于"功能-结构-工序-资源"多维本体的
装配资源知识图谱（ARKG），通过本体映射将CAD模型的实例化数据融入知识图谱，
实现了产品、工序和装配资源信息的有效耦合；
但该工作主要面向CAD属性数据，未涉及几何约束的系统性解析。

事实上，STEP/ISO 10303\cite{srinivasan2017bom}格式的CAD文件以标准化方式编码了
零件的父子包含关系、几何配合约束（同轴、贴合、平行等）和制造特征（PMI）——
这些信息以确定性方式存在，精度高且无需人工标注。
本文提出的CAD解析轨道正是利用OCCT从STEP文件中
提取装配层级树（\textit{isPartOf}）、配合约束（\textit{matesWith}）
与空间邻接关系（\textit{adjacentTo}），
弥补了纯文本知识图谱在几何拓扑知识方面的空白，
构成本文相较于现有工业 RAG 工作的主要差异之一。

\subsubsection{BOM驱动的实体注册与对齐}

BOM（Bill of Materials）是制造业中零件身份的权威来源，
包含零件编号（PN）、名称、数量和材料规格等结构化属性\cite{srinivasan2017bom}。
Liu等人\cite{liu2023avkg}指出，
航空装配领域的数据源高度非结构化，实体边界模糊
（如"雷达校准"可归类为操作实体，也可将"雷达"单独归为部件实体），
这一问题在跨源实体对齐时尤为突出。
以BOM作为图谱实体节点的注册基础，可有效解决多源数据中零件名称不一致的问题：
手册中的"主轴承"经实体链接与BOM中的 PN:~1002-AX 对应，
CAD中的几何体引用亦通过STEP文件的\texttt{PRODUCT}实体表解析后与BOM对齐，
从而在三源异构数据之间建立统一的实体标识空间。

\subsubsection{多智能体知识抽取}

由于三类数据源的格式与语义差异显著，单一流水线难以高质量处理全部来源。
Guo等人\cite{guo2024multiagent}证明，多智能体系统（Multi-Agent System）
通过分工协同——各智能体专注于自身最擅长的数据类型——总体性能优于单体方案。
Park等人\cite{park2023generative}进一步展示了多智能体协作在复杂推理任务中的潜力。
受此启发，本文设计了BOM解析智能体、文本抽取智能体、
CAD结构化解析智能体与验证智能体等专职智能体，
协同完成三源知识的抽取、融合与一致性校验。

\subsubsection{图谱推理与RAG融合}

GraphRAG\cite{edge2024graphrag}通过识别命名实体为锚节点、可变深度邻域扩展完成多跳遍历，
有效支持"某零件的所有子装配体"等结构化查询，
但其图谱仅来源于文本抽取，不含几何信息；
此外，传统 GraphRAG 多依赖关键词字符串匹配定位锚节点，
对中英混合查询、近义术语场景的鲁棒性有限。
本文的知识图谱检索路径（路径三）在此基础上做出两项改进：
其一，以CAD解析所得的几何约束节点丰富了图谱的结构深度，
从而支持纯文本GraphRAG无法回答的几何拓扑类查询；
其二，将锚节点定位从关键词CONTAINS匹配改为基于BGE-M3向量的密集语义近邻搜索，
使中英混合查询在统一编码空间内直接命中相关实体，再以一跳邻域扩展形成上下文子图。
这一设计减少了对精确字符串命中的依赖，为跨章节、跨语言别名查询提供了更稳定的锚节点定位方式。

% ─────────────────────────────────────────────────────────────────────────────
\subsection{查询扩展与生成增强}
\label{sec:related-generation}

\subsubsection{多查询扩展}

多查询扩展（Multi-Query Expansion）是缓解词汇不匹配的常用技术。
Ma等人\cite{ma2023query}证明，
利用LLM将原始查询改写为多种不同表述后并行检索、
再对结果合并去重，可显著提升召回覆盖率。
本文将多查询扩展应用于全部三条检索路径，
构成查询集~$\mathcal{Q} = \{q, q_1, q_2\}$，
是本系统实现高召回率的重要技术手段之一。

% ─────────────────────────────────────────────────────────────────────────────
\subsection{航空发动机装配智能化}
\label{sec:related-aero}

航空发动机装配是制造业中精度要求极高、工艺高度复杂的环节之一。
传统装配辅助系统多依赖人工编写的结构化数据库和固定检索规则，
缺乏对自然语言查询的理解能力，在面对机型更新或工艺变更时适应性较差。
近年来，随着知识图谱与大语言模型技术的快速发展，
该领域的智能化研究日趋活跃，可从以下三个维度加以梳理。

\textbf{航空领域知识图谱应用。}
Tang等人\cite{tang2023aircraft}系统综述了航空故障诊断知识图谱的构建与应用研究；
Meng等人\cite{meng2024aircraft}构建了面向航空PHM（预测与健康管理）的故障知识图谱平台，
涵盖航空发动机典型故障的实体建模和推理路径；
Wu等人\cite{wu2024engine}专门构建了航空发动机润滑系统故障知识图谱，
并在实际工程场景中验证了KG辅助故障诊断的有效性。
上述工作一致表明，专业化知识图谱是提升航空领域故障诊断精度的关键手段；
然而它们均聚焦于故障诊断场景，
未涉及装配引导这一需要三源异构知识深度融合的精细任务。

\textbf{航空装配知识图谱构建。}
Liu等人\cite{liu2023avkg}利用联合知识抽取模型自动从装配手册文本中
提取实体和关系，构建了含5类实体（步骤、操作、工具、部件、设施）和
3类关系（Component-Whole、Content-In、Instrument-Agency）的航空装配KG，
以Neo4j作为图数据库存储，实体和关系抽取的F1-score分别达89.71\%和91.27\%；
Shi等人\cite{shi2022arkg}构建了基于"功能-结构-工序-装配资源"多维本体的
装配资源知识图谱，将CAD实例数据和工程知识在统一语义框架下集成，
为装配工艺设计提供知识共享与复用支持。
两项工作均验证了KG在装配领域的可行性，
但其数据融合粒度仍有限，
尚未将BOM身份注册、技术手册工序知识与三维几何约束纳入统一图谱框架。

\textbf{LLM与装配智能化融合。}
Blumenthal等人\cite{blumenthal2020nlp}探索了将技术手册文本结构化以支持关键字检索的方案；
Zhao等人\cite{zhao2022assembly}提出基于BERT\cite{devlin2019bert}的装配指令理解模型，
在有监督条件下提升了指令语义解析能力。
然而，上述工作均未引入大语言模型的生成能力，难以应对表述多样的开放域问题。
Zhao等人\cite{zhao2023llmsurvey}的综述表明，LLM在自然语言理解、
推理和生成方面具有显著的涌现能力；
Wu 等\cite{wu2024engine}将 LLM 与航空发动机润滑系统 KG 协同，
证明了 LLM 在装备级专业故障问答中的应用潜力，
但纯LLM方案存在幻觉和知识陈旧问题，
在涉及精确数值（扭矩、公差、零件编号）的查询上可靠性不足\cite{fan2024ragsurvey}。

Liu等人\cite{liu2024jointkg}则进一步将航空装配知识图谱通过前缀微调嵌入LLM，
利用2-step子图驱动前缀调优，在功能测试场景下实现了98.5\%的故障诊断准确率，
验证了KG结构化知识与LLM推理能力深度融合的可行性。

\textbf{与本文最相关的工作。}
Shen等人\cite{shen2025hybridrag}的HybridRAG
在飞机故障排查场景下证明了KG+向量混合检索的有效性（F1提升4\%、幻觉率降低7\%），
并设计了Planner-Agent迭代探索算法，通过计划-执行-自我反思的循环提升答案质量。
本文借鉴其KG增强检索思想，并针对航空发动机装配的特殊需求做出三项关键扩展：
（1）引入三源（BOM+CAD+手册）多智能体知识图谱构建流水线，
将本体规模从手册文本的5类实体扩展至7类实体和9类关系，
构建覆盖更完整的装配知识底座；
（2）将KG锚节点定位从关键词匹配改为BGE-M3语义近邻搜索，
提升中英混合查询与专业别名查询下的图谱召回稳定性；
（3）将线性加权融合改进为RRF无参融合，提升文本块与图谱子图融合的鲁棒性。
上述扩展使本系统在知识覆盖度、结构化推理与融合可靠性上
相对现有工作具有差异化优势，并直接面向装配引导这一高精度应用场景。



% ═══════════════════════════════════════════════════════════════════════════
% 第 3 章  方法
% ═══════════════════════════════════════════════════════════════════════════
% ═════════════════════════════════════════════════════════════════════════════
%  第三章  方法
% ═════════════════════════════════════════════════════════════════════════════
\section{方法}
\label{sec:method}

% ─────────────────────────────────────────────────────────────────────────────
\subsection{系统总体架构}
\label{sec:overview}

本文提出的面向航空发动机装配的多模态混合检索增强生成框架整体分为
\textbf{离线知识构建}和\textbf{在线问答推理}两个阶段，系统架构如图~\ref{fig:arch}~所示。

\textbf{离线阶段}分两条并行流水线：
第一条处理非结构化的装配技术手册，经文档解析、语义切块、多模态向量化后存入向量数据库；
第二条处理结构化的CAD模型与手册文本，通过双轨知识图谱构建方法将零件结构关系、
装配工序、配合约束等知识写入图数据库。
两类知识源相互补充，共同构成系统的统一知识底座。

\textbf{在线阶段}：用户输入自然语言问题后，系统首先通过大语言模型（LLM）进行
多查询扩展，再并行启动四条检索路径——密集向量检索、BM25稀疏检索、
跨模态图像检索、知识图谱检索——通过互惠排名融合（RRF）将四路结果合并为
统一排名，经交叉编码器精排后送入LLM完成答案生成，
最终以服务器推送事件（SSE）流式返回给用户。

% 图占位符（替换为实际图片时取消注释）
% \begin{figure}[htbp]
%   \centering
%   \includegraphics[width=\linewidth]{figures/architecture.pdf}
%   \caption{系统总体架构图}
%   \label{fig:arch}
% \end{figure}

% ─────────────────────────────────────────────────────────────────────────────
\subsection{双轨知识图谱构建}
\label{sec:kg}

航空发动机装配领域存在两类本质不同的知识来源：
编码在CAD三维模型中的精确结构化数据，以及散布在装配技术手册文本中的工艺程序性知识。
本文为两类来源设计了独立的抽取流水线，最终汇入统一的Neo4j图数据库，
形成完整的装配知识图谱~$\mathcal{G} = (\mathcal{V}, \mathcal{E})$。

% ·············································································
\subsubsection{知识图谱本体设计}
\label{sec:ontology}

为规范图谱的语义边界，本文结合航空发动机装配领域特点，
设计了包含6类实体和7类关系的本体模式，如表~\ref{tab:ontology}~所示。

\begin{table}[htbp]
  \centering
  \caption{航空发动机装配知识图谱本体模式}
  \label{tab:ontology}
  \renewcommand{\arraystretch}{1.3}
  \begin{tabularx}{\linewidth}{llX}
    \toprule
    \textbf{类别} & \textbf{名称} & \textbf{说明} \\
    \midrule
    \multirow{6}{*}{\textbf{实体}}
      & Part（零件）       & 具有明确名称的单体零部件，如"高压压气机第3级叶片" \\
      & Assembly（子系统） & 包含子零件的装配体，如"高压压气机" \\
      & Procedure（工序）  & 具体装配操作动作，如"施加预紧扭矩" \\
      & Tool（工具）       & 装配工具器具，如"T-205专用锁紧工具" \\
      & Specification（规范）& 含具体数值和单位的技术要求，如"50\,N$\cdot$m"、"0.05\textasciitilde0.12\,mm" \\
      & Interface（接口）  & 零件间的配合面，如"榫头-榫槽配合" \\
    \midrule
    \multirow{7}{*}{\textbf{关系}}
      & precedes          & 工序A必须在工序B之前完成（有向序） \\
      & participatesIn    & 零件/子系统参与某工序 \\
      & requires          & 工序需要某工具 \\
      & specifiedBy       & 工序/接口受某规范约束 \\
      & matesWith         & 零件间存在配合关系 \\
      & hasInterface      & 零件拥有某配合接口 \\
      & constrainedBy     & 工序/零件受几何约束限制 \\
    \bottomrule
  \end{tabularx}
\end{table}

本体模式在知识抽取时作为强约束注入提示词，确保图谱边界清晰、无类型歧义。

% ·············································································
\subsubsection{CAD结构化解析（轨道一）}
\label{sec:cad}

给定描述发动机装配体的STEP格式CAD文件~$\mathcal{F}$，
本文使用Open CASCADE Technology（OCCT）的Python绑定\texttt{pythonocc-core}
从中提取三类结构化知识。

\paragraph{（1）装配层级提取}

STEP文件中的\texttt{NEXT\_ASSEMBLY\_USAGE\_OCCURRENCE}（NAUO）记录
编码了零件的父子包含关系。
对文件中每一条NAUO记录~$(r_p, r_c)$，
通过\texttt{PRODUCT}实体表将引用编号解析为人类可读的零件名称，
并生成\texttt{isPartOf}有向三元组：

\begin{equation}
  \langle \text{child},\ \texttt{isPartOf},\ \text{parent} \rangle
  \label{eq:ispartof}
\end{equation}

全部三元组构成完整的装配层级树~$\mathcal{T} \subseteq \mathcal{G}$，
直接反映发动机从顶层总装体到最底层零件的完整包含结构。

\paragraph{（2）配合约束提取}

在NAUO关联关系和\texttt{GEOMETRIC\_TOLERANCE}记录的基础上，
提取零件间的装配配合关系，生成带属性的三元组：

\begin{equation}
  \langle \text{part}_a,\ \texttt{matesWith},\ \text{part}_b \rangle
  \quad \text{含属性}\ \{\mathit{constraint\_type},\ \mathit{interface}\}
  \label{eq:mateswith}
\end{equation}

\paragraph{（3）空间邻接关系提取}

对STEP文件中每对传递后的实体几何体~$(\sigma_i, \sigma_j)$，
计算其轴对齐包围盒（AABB）$\mathcal{B}_i$ 和 $\mathcal{B}_j$ 之间的间距：

\begin{equation}
  \delta_{ij} =
  \sqrt{\sum_{u \in \{x,y,z\}}
    \max\!\left(0,\ \max\bigl(u_i^{\min}, u_j^{\min}\bigr)
                     - \min\bigl(u_i^{\max}, u_j^{\max}\bigr)\right)^{\!2}}
  \label{eq:gap}
\end{equation}

当~$\delta_{ij} < \tau$（默认阈值~$\tau = 2\,\text{mm}$）时，
判定两零件在物理上相邻，写入：

\begin{equation}
  \langle \text{part}_a,\ \texttt{adjacentTo},\ \text{part}_b \rangle
  \quad \text{含属性}\ \{gap\_mm = \delta_{ij}\}
  \label{eq:adjacent}
\end{equation}

与依赖标注语料的NLP三元组抽取方法相比，
CAD直接解析方法从几何数据中确定性地获取零件结构关系，
精度高且无需任何训练数据，是本文在知识获取方式上的核心创新之一。

% ·············································································
\subsubsection{LLM零样本三元组抽取（轨道二）}
\label{sec:llm-kg}

CAD几何数据仅能覆盖结构性知识，而装配力矩要求、工序顺序约束、
工具选用规范等程序性知识只存在于技术手册文本中。
本文采用LLM零样本抽取流水线，从手册文本块中自动获取这类知识。

对文档解析后的每个文本块~$c_i$，
构造携带完整本体模式的结构化提示词，
驱动LLM输出JSON格式的实体-关系集合，形式化定义为：

\begin{equation}
  \mathcal{T}_i = \Phi_{\text{LLM}}(c_i,\ \mathcal{O})
                = \bigl\{(e_h^{(k)},\ r^{(k)},\ e_t^{(k)})\bigr\}_{k=1}^{K_i}
  \label{eq:llm-extract}
\end{equation}

其中~$\mathcal{O}$~为本体模式，$e_h$、$e_t$~分别为头实体和尾实体，
$r$~为关系类型。
提示词中内置了三类Few-shot示例（工序链、配合关系、工具规范），
并通过\textbf{二轮Gleaning机制}对首轮结果进行补全——
第二轮提示词明确要求LLM只输出首轮遗漏的实体和关系，有效减少知识遗漏。

提取完成后，通过\textbf{三级级联实体对齐}消除同义表达，保证图谱的一致性：
\begin{enumerate}
  \item \textbf{规则级}：基于预置的航空领域缩写词典
        （如 HPC$\rightarrow$高压压气机、HPT$\rightarrow$高压涡轮）进行精确映射；
  \item \textbf{模糊级}：使用编辑距离（SequenceMatcher）
        对相似度高于阈值的实体名称进行合并；
  \item \textbf{语义级}：对前两级未能对齐的候选实体，
        使用BGE-M3语义向量余弦相似度作最终判断。
\end{enumerate}

此外，针对\texttt{precedes}关系可能引入的工序环路，
采用Kahn拓扑排序算法进行有向无环图（DAG）验证，
一旦检出环路则截断置信度最低的边，保证装配工序图的有效性。

两轨数据最终通过\texttt{MERGE}操作叠加写入同一Neo4j实例，
构成统一的装配知识图谱~$\mathcal{G}$。

% ─────────────────────────────────────────────────────────────────────────────
\subsection{文档解析与多模态向量化}
\label{sec:index}

\subsubsection{文档解析流水线}
\label{sec:parse}

航空发动机装配手册通常为图文混排的PDF文件，本文采用双路解析策略：
\begin{itemize}
  \item \textbf{可复制PDF}：使用PyMuPDF直接提取文本层，保留原始字符坐标；
  \item \textbf{扫描版/复杂版式PDF}：路由至DeepDoc AI解析引擎，
        执行版面检测（Layout Analysis）和OCR，
        并按ATA章节编号进行结构化分块。
\end{itemize}

文本按句子边界切块，目标块长度为500字符。
页面中尺寸大于~$100 \times 100$~像素的嵌入图片被提取并持久化存储，
供后续多模态索引使用。

\subsubsection{多模态向量化与双索引构建}
\label{sec:vectorize}

每个文本块~$c_i$~经\textbf{BGE-M3}编码器
（输出维度~$d_{\text{text}} = 1024$）得到密集向量：

\begin{equation}
  \mathbf{v}_i^{\text{text}} = \text{BGE-M3}(c_i) \in \mathbb{R}^{1024}
  \label{eq:text-vec}
\end{equation}

对每张提取图片~$f_j$，首先调用视觉大模型（MiniMax M2.5 Vision）
生成中文技术描述~$\hat{c}_j$（图注），
再经BGE-M3编码得到文字代理向量~$\mathbf{v}_j^{\text{cap}}$，
用于支持"文字查图"路径。
同时，图片原图经\textbf{Chinese-CLIP}
（输出维度~$d_{\text{img}} = 768$）编码得到跨模态图像向量：

\begin{equation}
  \mathbf{v}_j^{\text{img}} = \text{CLIP}_{\text{image}}(f_j) \in \mathbb{R}^{768}
  \label{eq:img-vec}
\end{equation}

全部向量存入Qdrant向量数据库，在同一集合下以具名向量字段
\texttt{text\_vec}和\texttt{image\_vec}分别存储，
支持两个向量空间的独立检索。

% ─────────────────────────────────────────────────────────────────────────────
\subsection{多维混合检索}
\label{sec:retrieval}

本节是本文的核心贡献。
给定用户查询~$q$，系统并行执行四条检索路径并融合结果。

\subsubsection{多查询扩展}
\label{sec:multi-query}

为缓解词汇表不匹配问题、扩大召回覆盖面，
首先用LLM将原始查询改写为两个不同表述角度，构成查询集：

\begin{equation}
  \mathcal{Q} = \{q,\ q_1,\ q_2\}
              = \{q\} \cup \Phi_{\text{LLM}}^{\text{expand}}(q)
  \label{eq:multi-query}
\end{equation}

后续四条检索路径均在~$\mathcal{Q}$~中的每个查询上独立执行，
最终合并去重。

\subsubsection{密集向量检索（路径一）}
\label{sec:dense}

对查询集中每个查询~$q' \in \mathcal{Q}$，
用BGE-M3编码得到1024维查询向量~$\mathbf{q}^{\text{text}}$，
在Qdrant的\texttt{text\_vec}命名向量空间中通过HNSW图近似最近邻搜索，
按余弦相似度排序返回Top-$N$候选：

\begin{equation}
  \text{score}_{\text{dense}}(q', c_i)
  = \frac{\mathbf{q}^{\text{text}} \cdot \mathbf{v}_i^{\text{text}}}
         {\|\mathbf{q}^{\text{text}}\|\ \|\mathbf{v}_i^{\text{text}}\|}
  \label{eq:dense-score}
\end{equation}

密集检索擅长捕捉语义相似性，适合处理表述与手册用词不同但语义等价的查询。

\subsubsection{BM25稀疏检索（路径二）}
\label{sec:bm25}

装配领域中，零件编号（如"T-205"）、公差数值（如"0.05\textasciitilde0.12\,mm"）、
扭矩规格等精确标识符对检索准确性至关重要，
而密集向量检索往往因词汇平滑而弱化这类精确匹配能力。
为此，本文引入BM25稀疏检索作为补充路径。

考虑到手册中中英文混排的特点，本文设计\textbf{双轨分词}策略：
中文词汇通过jieba进行语义分词；
零件编号、测量数值等字母数字混合串通过正则表达式分词器处理，
避免被切断。
对查询~$q'$~和文档~$d$，BM25评分公式为：

\begin{equation}
  \text{score}_{\text{BM25}}(d, q')
  = \sum_{i=1}^{M} \text{IDF}(q'_i)
    \cdot \frac{f(q'_i,\, d) \cdot (k_1 + 1)}
               {f(q'_i,\, d)
                + k_1\!\left(1 - b + b\,\dfrac{|d|}{\text{avgdl}}\right)}
  \label{eq:bm25}
\end{equation}

其中

\begin{equation}
  \text{IDF}(q'_i)
  = \log\!\left(\frac{N - n(q'_i) + 0.5}{n(q'_i) + 0.5} + 1\right)
  \label{eq:idf}
\end{equation}

$N$~为语料总文档数，$n(q'_i)$~为包含词项~$q'_i$~的文档数。
超参数设置为~$k_1 = 1.5$，$b = 0.75$。
BM25侧的召回数设为密集检索的2倍（$2N$），
以充分利用其计算代价低廉的优势扩大候选池。

\subsubsection{跨模态图像检索（路径三）}
\label{sec:clip}

航空发动机装配手册中包含大量装配截面图、爆炸视图、公差标注图，
这些图纸承载着纯文本无法完整表达的装配知识。
这是本文相较于现有RAG方法的重要区别。

对查询~$q'$，使用Chinese-CLIP的文本编码器生成768维跨模态查询向量
$\mathbf{q}^{\text{clip}}$，在Qdrant的\texttt{image\_vec}命名向量空间中检索，
并通过元数据过滤器限定只返回\texttt{chunk\_type=image}的块：

\begin{equation}
  \text{score}_{\text{CLIP}}(q', f_j)
  = \frac{\mathbf{q}^{\text{clip}} \cdot \mathbf{v}_j^{\text{img}}}
         {\|\mathbf{q}^{\text{clip}}\|\ \|\mathbf{v}_j^{\text{img}}\|}
  \label{eq:clip-score}
\end{equation}

每个查询返回Top-$\lfloor N/2 \rfloor$~张图片，
连同图片的URL地址和LLM生成的中文图注共同作为多模态上下文传入后续流程。
这使得系统能够在最终回答中直接引用装配图纸，
显著提升图纸相关问题的可解释性。

\subsubsection{知识图谱检索（路径四）}
\label{sec:kg-retrieval}

对涉及零件组成关系、装配约束、工序依赖等结构化推理的查询，
系统向Neo4j知识图谱~$\mathcal{G}$~发起Cypher多跳查询。
形式化地，目标是找到最能支撑回答~$q$~的子图：

\begin{equation}
  \mathcal{G}^{*}
  = \arg\max_{\mathcal{G}' \subseteq \mathcal{G}}\;
    P\!\left(\mathcal{G}' \mid q,\; \mathcal{G}\right)
  \label{eq:kg-retrieval}
\end{equation}

实现上，从查询~$q$~中识别命名实体作为锚节点，
通过可变深度邻域扩展进行多跳遍历。
以下给出两类典型查询模式示例：

\begin{lstlisting}[language=SQL, caption={Cypher多跳查询示例}]
-- 查询零件的全部子装配体（深度3跳以内）
MATCH (root:Part {part_name: $name})<-[:isPartOf*1..3]-(child)
RETURN child

-- 查询某工序所需全部工具及其规范
MATCH (p:Procedure {text: $proc})-[:requires]->(t:Tool)
OPTIONAL MATCH (p)-[:specifiedBy]->(s:Specification)
RETURN t, s
\end{lstlisting}

检索得到的子图被序列化为结构化文本（节点属性 + 边标签），
纳入后续检索结果池参与融合。

\subsubsection{四路RRF融合}
\label{sec:rrf}

四条检索路径返回的候选集处于各自独立的评分空间，无法直接比较。
本文通过互惠排名融合（Reciprocal Rank Fusion，RRF）将其统一到同一排名空间。
设~$\mathcal{R} = \{\text{dense},\ \text{BM25},\ \text{CLIP},\ \text{KG}\}$~为四个检索器，
对候选块~$d$~的RRF综合得分为：

\begin{equation}
  \text{RRF}(d)
  = \sum_{r \in \mathcal{R}} \frac{1}{k + \text{rank}_r(d)}
  \label{eq:rrf}
\end{equation}

其中~$k = 60$~为平滑常数，用于抑制任意单路高排名带来的过度支配；
若候选~$d$~未出现在检索器~$r$~的结果中，则该项贡献为0。
全部候选按RRF得分降序排列，取前~$L$~个传入精排阶段。

RRF的核心优势在于：无需对各路分数进行归一化，
对单路召回失败具有鲁棒性，
且可自然地将文本块、图片块、图谱子图三种异质内容统一排序。

% ─────────────────────────────────────────────────────────────────────────────
\subsection{交叉编码器精排}
\label{sec:rerank}

RRF融合的粗排结果基于各路独立编码，
未能充分利用查询与候选段落之间的细粒度交互信息。
为此，本文引入交叉编码器（Cross-Encoder）对粗排结果进行精排，
使用BGE-Reranker-Base作为精排模型。

给定查询~$q$~和候选段落~$d$，将两者拼接为如下格式输入序列：

\begin{equation}
  \text{Input} = [\text{CLS}]\; q\; [\text{SEP}]\; d\; [\text{SEP}]
  \label{eq:rerank-input}
\end{equation}

该序列经过多层Transformer编码：

\begin{equation}
  \mathbf{H} = \text{Transformer}(\text{Input})
  \label{eq:rerank-transformer}
\end{equation}

取[CLS]位置的隐状态~$\mathbf{h}_{\text{CLS}} = \mathbf{H}[\text{CLS}]$~
作为查询-段落对的联合语义表示，
经前馈网络输出相关性得分：

\begin{equation}
  \text{Score}(q, d) = \text{FFN}(\mathbf{h}_{\text{CLS}})
  \label{eq:rerank-score}
\end{equation}

与双编码器相比，交叉编码器在推理时对查询和文档进行全注意力交互，
能更精准地捕捉语义对齐关系。
按相关性得分降序取前~$K$（默认~$K=4$）个作为最终上下文。

% ─────────────────────────────────────────────────────────────────────────────
\subsection{答案生成}
\label{sec:generation}

将Top-$K$~上下文片段（含文本块、图片图注及对应图片URL）
按相关性得分依次拼接，构成结构化提示词送入LLM。生成目标为：

\begin{equation}
  \hat{a} = \arg\max_{a}\;
  \prod_{j=1}^{|a|}
  P\!\left(a_j \mid a_{<j},\; q,\; \mathcal{C}\right)
  \label{eq:generation}
\end{equation}

其中~$\mathcal{C} = \{d_1, \ldots, d_K\}$~为最终上下文集合，
$a_{<j}$~为已生成的前序词元。

\textbf{系统提示设计}方面，为支持交互式装配引导场景，
本文采用结构化指令约束LLM的输出行为：
助手必须首先确认当前装配状态，再给出下一步可操作指令，
并等待用户反馈后方可继续。
这一设计有效避免了未经核实直接给出多步操作的风险，
契合航空发动机装配对操作安全性的高要求。

答案以\textbf{服务器推送事件（SSE）}流式返回前端，由三类帧组成：
\texttt{stage}帧（当前检索/生成阶段提示）、
\texttt{delta}帧（LLM增量输出词元）、
\texttt{done}帧（来源引用列表和图片URL集合）。
   % ← 已完成，直接引入


% ═══════════════════════════════════════════════════════════════════════════
% 第 4 章  实验与结果
% ═══════════════════════════════════════════════════════════════════════════
\section{实验与结果}

\subsection{实验配置}
% TODO：硬件环境、模型版本、超参数列表

\subsection{数据集构建}
% TODO：语料来源、标注方案、统计信息表

\subsection{评价指标}
% TODO：Precision@K、Recall@K、F1、Hit Rate、MRR、端到端 ROUGE/BLEU

\subsection{基线方法}
% TODO：Naive RAG / BM25-only / Dense-only / KG-only / HybridRAG-原版

\subsection{主实验结果}
% TODO：主对比表（Table 2）+ 分析

\subsection{消融实验}
% TODO：逐路径消融（去掉 BM25 / CLIP / KG）+ Table 3

% 完成后取消注释：% ═════════════════════════════════════════════════════════════════════════════
%  第四章  实验与结果
% ═════════════════════════════════════════════════════════════════════════════
\section{实验与结果}
\label{sec:experiment}

% ─────────────────────────────────────────────────────────────────────────────
\subsection{实验配置}
\label{sec:exp-setup}

实验运行于一台 Windows 11 工作站，
硬件配置为 Intel Core i7-13 系列处理器、32\,GB DDR5 内存与
NVIDIA GeForce RTX 5070 GPU。
本地部署的模型组件包括 BGE-M3 文本编码器（1024 维）与
BGE-Reranker-Base 交叉编码器精排模型，均通过 PyTorch 加载至 GPU 进行推理。
向量数据库选用 Qdrant 本地文件模式，图数据库为 Neo4j 5 社区版（Docker 容器）。
大语言模型采用主备双线策略：主模型为 MiniMax M2.5（OpenAI 兼容协议），
备用模型为 SiliconFlow 平台部署的 Qwen2.5-14B-Instruct，
当主线返回 4xx 或超时时由 \texttt{FallbackLLMClient} 自动切换。

% ─────────────────────────────────────────────────────────────────────────────
\subsection{数据集构建}
\label{sec:exp-dataset}

\subsubsection{知识源语料}

实验所用语料来源于 Pratt \& Whitney Canada PT6A 系列航空发动机的真实工程文档，
覆盖整机 10 个 ATA 章节（61/70/71/72/73/74/75/76/77/79）。
文档总规模与三源 KG 构建产出统计如表~\ref{tab:dataset-stats}~所示。

\textbf{数据来源与合规性}：
PT6A 维修手册（Maintenance Manual P/N 3013242）、零件目录（IPC P/N 3013243）
及训练用 STEP 模型（C52696C / C89119）来自可公开获取的
航空培训与售后服务资料。
本文仅在学术研究环境中使用上述文档，未对外二次分发原始 PDF 文件；
评测问答集与三元组结构化结果可按本文实验配置导出，
以便在授权范围内复现实验。
所述数据的知识产权、使用许可与出口管制责任仍以原始版权方和授权条款为准。

\begin{table}[htbp]
  \centering
  \caption{PT6A 三源知识图谱构建统计}
  \label{tab:dataset-stats}
  \renewcommand{\arraystretch}{1.3}
  \begin{tabularx}{\linewidth}{lXr}
    \toprule
    \textbf{数据源} & \textbf{描述} & \textbf{规模} \\
    \midrule
    \multirow{4}{*}{\textbf{文本语料}}
      & PT6A 维修手册（textin OCR Markdown 三分册） & 1.78\,MB \\
      & 整机 ATA 章节引用 & 1100+ 次 \\
      & 文本块（chunk）数（BGE-M3 编码） & 3{,}376 \\
      & 整机表格（HTML 形式保留） & 198 \\
    \midrule
    \multirow{2}{*}{\textbf{结构化}}
      & BOM 表（IPC 风格 docx 转 csv） & 581 子表 / 3{,}717 PN \\
      & CAD STEP 文件（C52696C / C89119） & 2 个 \\
    \midrule
    \multirow{4}{*}{\textbf{KG 三元组}}
      & Stage 1（BOM）：isPartOf + interchangesWith & 3{,}586 \\
      & Stage 2（手册 LLM 抽取）：7 类关系 & 6{,}022 \\
      & Stage 3（CAD）：5 类关系 & 1{,}209 \\
      & 后处理 + 跨源对齐写入 Neo4j & \textbf{8{,}752 边 / 7{,}333 节点} \\
    \bottomrule
  \end{tabularx}
\end{table}

\subsubsection{评测问答集}

为兼顾评测严谨性与系统覆盖广度，本文构建了双层评测集：

\textbf{Tier-1 精评集（60 题）}~聚焦于压气机子系统（ATA 72-30），
来源为人工精读维修手册第 72-30 章后建立的黄金三元组集
（108 实体 / 99 三元组 / 9 类关系全覆盖），通过模板化反向生成
零件结构、工序工具、技术规范、几何配合四类问题，
分布严格按黄金集真实关系比例（A 20 / B 12 / C 9 / D 19 题）。
每题包含中英双语问题与参考答案、来源 ATA 章节、关联子图标识，
并通过 BGE-M3 在 Qdrant 中反向召回 Top-5 候选 chunks 作为评测锚点。
该层评测更适合验证系统对已知结构关系的召回与生成能力，
并不等同于真实生产环境中的开放式自然提问。

\textbf{Tier-2 粗评集（20 题）}~跨章节覆盖涡轮装配（72-50）、
燃油系统（73-10/20）、控制与指示（76-77）、滑油系统（79）四个子系统，
每题以高关联度零件作为锚点，利用 KG 子图生成装配场景问题。
该层评测不提供 \texttt{gold\_chunk} 标注，仅以 LLM-Judge 进行半自动评分，
用于验证系统在跨章节场景下的鲁棒性；其结论主要作为补充证据。

% ─────────────────────────────────────────────────────────────────────────────
\subsection{评价指标}
\label{sec:exp-metrics}

\subsubsection{检索阶段指标（Tier-1）}

对每条问题 $q$，给定系统返回的 Top-$K$ 候选块集合 $R_K$ 与黄金候选集 $G$，定义：

\begin{equation}
  \text{Recall@}K
  = \frac{|R_K \cap G|}{|G|}
  \label{eq:recall}
\end{equation}

\begin{equation}
  \text{MRR}
  = \frac{1}{|Q|} \sum_{q \in Q} \frac{1}{\mathrm{rank}_q^{*}}
  \label{eq:mrr}
\end{equation}

其中 $\mathrm{rank}_q^{*}$ 为问题 $q$ 中第一个相关结果在 $R_K$ 中的位次。
进一步报告 Hit@1、Hit@5 与 NDCG@5 以反映 Top-$K$ 内排序质量：

\begin{equation}
  \text{NDCG@}K
  = \frac{1}{|Q|}\sum_{q \in Q}
    \frac{\sum_{i=1}^{K}\dfrac{\mathbb{1}[d_i \in G]}{\log_2(i+1)}}
         {\sum_{i=1}^{\min(|G|,K)}\dfrac{1}{\log_2(i+1)}}
  \label{eq:ndcg}
\end{equation}

\subsubsection{生成阶段指标}

由于参考答案高度抽象（来自人工精读黄金集），
而系统生成答案多基于手册原文段落，二者在表述粒度上天然不一致，
本文采用 \textbf{LLM-as-Judge} 三维评分作为主要生成指标：

\begin{itemize}
  \item \textit{Relevance}（相关性）—— 答案与问题的相关程度（0--5）
  \item \textit{Completeness}（完整性）—— 答案是否覆盖了关键信息（0--5）
  \item \textit{Correctness}（正确性）—— 答案与参考答案的事实一致程度（0--5）
\end{itemize}

由 Qwen2.5-14B-Instruct 担任 Judge，温度设为 0.1 以保证评分稳定。
每条 QA 取三维度得分的平均值作为综合质量指标。
本文未对所有样本执行重复评分或专家盲评，因此 LLM-Judge 结论应理解为半自动评测结果。

% ─────────────────────────────────────────────────────────────────────────────
\subsection{基线方法}
\label{sec:exp-baselines}

由于现有工业 RAG 工作（特别是 HybridRAG\cite{shen2025hybridrag}）未开源、
且本文主要贡献聚焦于三源 KG 构建与多路检索融合而非检索算法替换，
本文采用\textbf{内部消融基线}以保证比较可复现性。
所有基线统一使用本文知识库（Qdrant 文本向量 + BM25 索引 + Neo4j 整机 KG），
通过路径开关 \texttt{path\_mask} 配置不同检索路径组合（表~\ref{tab:baselines}）。

\begin{table}[htbp]
  \centering
  \caption{内部消融基线配置}
  \label{tab:baselines}
  \renewcommand{\arraystretch}{1.3}
  \begin{tabularx}{\linewidth}{llXc}
    \toprule
    \textbf{ID} & \textbf{方法} & \textbf{启用路径} & \textbf{Reranker} \\
    \midrule
    B1 & Naive RAG (Dense-only)        & BGE-M3 单路       & 否 \\
    B2 & BM25-only                     & 稀疏单路           & 否 \\
    B3 & Hybrid (Dense + BM25 RRF)     & 文本双路           & 是 \\
    B4 & MHRAG (Dense + BM25 + KG)     & 文本双路 + KG 子图 & 是 \\
    \bottomrule
  \end{tabularx}
\end{table}

三路检索结果的合并策略与第~\ref{sec:rrf}~节描述一致：
密集（Dense）与稀疏（BM25）通过互惠排名融合（RRF, $k=60$）合并；
图谱子图（KG）作为结构化上下文直接加入候选池。
对于启用 Reranker 的配置，候选结果再经交叉编码器精排取 Top-5 作为生成阶段的上下文；
未启用 Reranker 的单路基线则直接取对应检索路径的 Top-5。

% ─────────────────────────────────────────────────────────────────────────────
\subsection{主实验结果}
\label{sec:exp-main}

\subsubsection{Tier-1 精评：检索阶段对比}

表~\ref{tab:exp-tier1-main}~汇总 4 个基线配置在 Tier-1 精评集（60 题）上的检索阶段指标。
所有配置均达到 Hit@5 = 1.000，说明 PT6A 维修手册第 72-30 章的语料覆盖密度较高。
B1 Dense-only 在 Hit@1 上达到 0.983，体现 BGE-M3 在专业语料上的强表征能力。
引入 KG 子图后，B4 的 Recall@5 低于 B3（0.697 vs 0.769），
原因是当前检索评价分母主要以黄金文本块为主，而 KG 子图作为异质上下文会占用 Top-5 位置。
因此，检索阶段指标需要与生成阶段指标结合解读。

\begin{table}[htbp]
  \centering
  \caption{Tier-1 精评：检索阶段基线对比（60 题平均）}
  \label{tab:exp-tier1-main}
  \renewcommand{\arraystretch}{1.25}
  \begin{tabularx}{\linewidth}{lYYYYY}
    \toprule
    \textbf{基线} & Recall@5 & MRR & Hit@1 & Hit@5 & NDCG@5 \\
    \midrule
    B1 Dense-only         & \textbf{0.778} & \textbf{0.992} & \textbf{0.983} & 1.000 & \textbf{0.987} \\
    B2 BM25-only          & 0.719 & 0.938 & 0.900 & 1.000 & 0.951 \\
    B3 Hybrid (RRF)       & 0.769 & 0.983 & 0.967 & 1.000 & 0.974 \\
    B4 MHRAG              & 0.697 & 0.926 & 0.867 & 1.000 & 0.951 \\
    \bottomrule
  \end{tabularx}
  \\[2pt]
  \footnotesize{注：B4 的 Recall@5 下降并不必然表示回答质量下降；
  该现象反映了文本块评价标准与 KG 子图异质上下文之间的张力。}
\end{table}

\subsubsection{Tier-1 精评：生成阶段对比}

表~\ref{tab:exp-tier1-judge}~报告 LLM-Judge 三维评分。
B4 MHRAG 在 Correctness 维度取得数值最高的得分（2.35），
相对 B1 Dense-only 提升 10.8\%，相对 B3 Hybrid 提升 14.6\%。
这说明 KG 子图对事实核验和关系补全有直接帮助。
在 Relevance 与 Completeness 维度上，B4 同样高于 B3，
但提升幅度相对标准差较小，因此本文将其解读为稳定趋势而非统计显著结论。

\begin{table}[htbp]
  \centering
  \caption{Tier-1 精评：生成阶段 LLM-Judge 三维评分（0--5）}
  \label{tab:exp-tier1-judge}
  \renewcommand{\arraystretch}{1.25}
  \begin{tabularx}{\linewidth}{lYYY}
    \toprule
    \textbf{基线} & Relevance ($\mu\pm\sigma$) & Completeness ($\mu\pm\sigma$) & Correctness ($\mu\pm\sigma$) \\
    \midrule
    B1 Dense-only         & $3.38 \pm 1.39$ & $2.88 \pm 1.34$ & $2.12 \pm 1.80$ \\
    B2 BM25-only          & $3.43 \pm 1.20$ & $2.78 \pm 1.15$ & $2.00 \pm 1.51$ \\
    B3 Hybrid (RRF)       & $3.48 \pm 1.21$ & $2.88 \pm 1.32$ & $2.05 \pm 1.74$ \\
    B4 MHRAG              & $\textbf{3.50} \pm 1.27$ & $\textbf{3.02} \pm 1.27$ & $\textbf{2.35} \pm 1.68$ \\
    \bottomrule
  \end{tabularx}
  \\[2pt]
  \footnotesize{$\sigma$ 为 60 题样本标准差。Judge 三维内部 Pearson 相关性平均为 +0.848（n=12 维度对），详见附录~\ref{appendix:judge}。}
\end{table}

\begin{figure}[htbp]
  \centering
  \includegraphics[width=\linewidth]{figures/fig8_main_compare.pdf}
  \caption{Tier-1 主实验结果对比。左侧显示文本块检索指标，右侧显示生成质量指标；B4 在纯检索排序上并不占优，但在答案正确性与完整性上更稳定。}
  \label{fig:main-compare}
\end{figure}

\subsubsection{题型分组分析}

题型分组结果显示，KG 路径的收益并非均匀分布：
其主要价值集中在需要显式结构关系的问题上。
B 工序工具类问题依赖工序、工具和规范之间的连接关系，
KG 子图能直接提供 \textit{requires} 与 \textit{specifiedBy} 等边，
因此相较 Dense-only 有更明显收益。
C 技术规范类问题则受益于 KG 中对扭矩、间隙和材料要求的结构化聚合。
相反，D 几何配合类问题的答案锚点多出现在连续手册文本中，
密集-稀疏文本融合已经可以覆盖大部分信息，KG 的边际收益较小。

\subsubsection{Tier-2 跨章节粗评}

表~\ref{tab:exp-tier2}~报告 Tier-2 粗评（20 题，跨 72-50/73/76-77/79 四章节）的 LLM-Judge 评分。
该评测集无 \texttt{gold\_chunk\_ids} 标注，故仅报告生成阶段指标。

\begin{table}[htbp]
  \centering
  \caption{Tier-2 跨章节粗评：LLM-Judge 三维评分（20 题）}
  \label{tab:exp-tier2}
  \renewcommand{\arraystretch}{1.25}
  \begin{tabularx}{\linewidth}{lYYY}
    \toprule
    \textbf{基线} & Relevance & Completeness & Correctness \\
    \midrule
    B1 Dense-only         & 3.10 & 2.40 & 2.00 \\
    B2 BM25-only          & 3.65 & 2.80 & 2.65 \\
    B3 Hybrid (RRF)       & 3.25 & 2.55 & 2.05 \\
    B4 MHRAG              & \textbf{3.75} & \textbf{3.05} & \textbf{2.70} \\
    \bottomrule
  \end{tabularx}
\end{table}

Tier-2 跨章节场景下，B4 MHRAG 在三项生成指标上均取得数值最高值。
相对 BM25-only，B4 在 Relevance 上提升 2.7\%，在 Completeness 上提升 8.9\%，
在 Correctness 上提升 1.9\%。
这一结果表明，BM25 对跨章节问题中的精确术语仍然非常重要，
而 KG 子图能够进一步补足单一关键词检索难以覆盖的装配层级和关系链。

% ─────────────────────────────────────────────────────────────────────────────
\subsection{消融实验}
\label{sec:exp-ablation}

在当前三路框架下，本文重点比较三条核心路径的作用：
Dense 负责语义召回，BM25 负责精确字符串匹配，KG 负责结构化关系补全。
由表~\ref{tab:exp-tier1-main} 和表~\ref{tab:exp-tier1-judge} 可见，
单独依赖 Dense 可以取得更高的检索排序指标，但生成阶段 Correctness 低于 B4；
单独依赖 BM25 在跨章节粗评中具有较强 Relevance，但 Tier-1 Completeness 最低；
Dense+BM25 的混合检索提高了文本召回稳定性，
而进一步加入 KG 子图后，生成阶段 Correctness 从 B3 的 2.05 提升至 2.35。

因此，本文的核心消融结论是：
\textbf{文本双路检索保证覆盖，KG 子图提高事实约束。}
这种组合在纯检索指标上并不总是占优，
但在需要装配层级、工序依赖和技术规范的专业问答中，
能够提供更可靠的结构化上下文。

\begin{figure}[htbp]
  \centering
  \includegraphics[width=\linewidth]{figures/fig9_ablation_delta.pdf}
  \caption{三路检索消融视角下各基线相对 B4 MHRAG 的指标差异。文本单路或双路在部分检索指标上更高，但在 LLM-Judge 生成质量上整体低于 B4。}
  \label{fig:ablation-delta}
\end{figure}

% ─────────────────────────────────────────────────────────────────────────────
\subsection{工程性能与幻觉率分析}
\label{sec:exp-engineering}

参照 Shen 等\cite{shen2025hybridrag}在 Scientific Reports 发表的 HybridRAG
工程指标体系，本节从端到端延迟（Latency）、答非所问率（Hallucination Rate, HR）
与显存占用（VRAM）三个工程维度进行对比。

\textbf{Latency 计量方式}：单题平均延迟由检索阶段总耗时与生成阶段总耗时
分别除以题量再相加得到，等价于工程部署中单题端到端响应时间。
\textbf{HR 定义}：参照\cite{shen2025hybridrag}的"答非所问率"概念，
本文以 LLM-Judge Correctness 维度评分小于 2.0 视为答非所问，
统计该比例作为 HR 指标。
\textbf{VRAM 估算}：本文 LLM 通过 OpenAI 兼容 API 调用远端 vLLM 服务，
本地仅承担 BGE-M3 编码与 BGE-Reranker 精排推理，
表中显存为各模型 fp16 推理理论值与 CUDA 工作区开销之和。

\begin{table}[htbp]
  \centering
  \caption{各基线工程性能与幻觉率对比（80 题双层评测集）}
  \label{tab:exp-engineering}
  \small
  \begin{tabular}{lrrrrr}
    \toprule
    基线 & 检索 (s) $\downarrow$ & 生成 (s) $\downarrow$ & E2E (s) $\downarrow$ & HR (\%) $\downarrow$ & VRAM (GB) $\downarrow$ \\
    \midrule
    B1 Dense-only      & 2.17 & 3.05 & 5.22 & 46.25 & 1.60 \\
    B2 BM25-only       & \textbf{2.03} & \textbf{2.94} & \textbf{4.98} & \textbf{35.00} & --   \\
    B3 Hybrid (RRF)    & 3.82 & 3.40 & 7.21 & 43.75 & 1.60 \\
    B4 MHRAG           & 4.78 & 2.98 & 7.76 & \textbf{35.00} & 1.60 \\
    \midrule
    \multicolumn{2}{l}{HybridRAG\cite{shen2025hybridrag}（参考）} & --
    & 14.19 & 15.01 & 26.12 \\
    \bottomrule
  \end{tabular}
\end{table}

\begin{figure}[htbp]
  \centering
  \includegraphics[width=\linewidth]{figures/fig12_engineering.pdf}
  \caption{工程性能对比。KG 路径增加检索阶段耗时，但 B4 在答非所问率与正确性之间取得更稳妥的折中。}
  \label{fig:engineering}
\end{figure}

B4 MHRAG 的检索阶段延迟高于文本检索基线，
主要原因是 KG 路径需要先进行语义锚节点定位，再访问 Neo4j 扩展一跳邻域。
但其 HR 与 BM25-only 持平，并低于 Dense-only，
同时在 Correctness 上高于所有内部基线。
这说明 KG 路径的主要价值并非降低答非所问率的绝对比例，
而是在可接受延迟内增强答案的事实约束和结构化完整性。

与 HybridRAG\cite{shen2025hybridrag} 的工程指标只能作部署形态参考：
参考论文采用本地化 LLM 推理，需将完整模型权重常驻显存；
本文采用远端 vLLM 服务 + OpenAI 兼容 API 调用模式，
本地仅承担 embedding 与精排推理。
因此，本文的本地显存更低并不直接等价于算法层面的显存优势，
而是说明该架构更适合资源受限的工业现场部署。

综上，本文 MHRAG 的实验故事并不是"所有指标同时占优"，
而是在文本检索召回已经较强的基础上，
通过三源知识图谱把装配层级、工序依赖和技术规范显式注入上下文，
以适度的检索延迟换取更好的事实一致性与跨章节完整性。



% ═══════════════════════════════════════════════════════════════════════════
% 第 5 章  讨论
% ═══════════════════════════════════════════════════════════════════════════
\section{讨论}

\subsection{各检索路径贡献分析}
% TODO

\subsection{知识图谱质量对生成效果的影响}
% TODO

\subsection{局限性与未来工作}
% TODO

% 完成后取消注释：% ═════════════════════════════════════════════════════════════════════════════
%  第五章  讨论
% ═════════════════════════════════════════════════════════════════════════════
\section{讨论}
\label{sec:discussion}

% ─────────────────────────────────────────────────────────────────────────────
\subsection{各检索路径贡献分析}
\label{sec:disc-path}

三路并行检索架构（Dense / BM25 / KG）的核心假设是：
不同路径分别覆盖装配问答中的不同信息需求。
基于第~\ref{sec:exp-main}~节与~\ref{sec:exp-ablation}~节的实验数据，
本文将其作用概括为"语义召回、精确匹配、结构约束"三类互补能力。

\textbf{Dense 与 BM25 的互补性。}~密集向量检索（B1 Dense-only）擅长捕捉
"压气机前支撑轴"这类语义近邻查询，
在 Tier-1 精评的 Hit@1 上达到 0.983，体现出较强的首位命中能力。
BM25（B2）则在含精确零件号（如 \texttt{8210-210A}）、
扭矩数值（\texttt{32-36\,lb.in}）的字符串匹配场景上具有天然优势。
RRF 融合（B3 Hybrid）保留二者优势，
使系统既能处理同义改写，也能保留工业文档中不可模糊化的编号与数值。

\textbf{KG 路径的结构化推理。}~B4 MHRAG 在 Tier-1 生成阶段取得
Relevance 3.50、Completeness 3.02 与 Correctness 2.35。
其中 Correctness 相对 B3 Hybrid 提升 14.6\%，
说明 KG 子图对事实核验和关系补全有直接帮助。
进一步观察题型分组，KG 路径在工序工具类与技术规范类问题上贡献最明显：
前者依赖 \textit{requires}、\textit{precedes} 等工序关系，
后者依赖 \textit{specifiedBy} 对扭矩、间隙和材料要求的结构化聚合。
这也是本文三源 KG 相对纯文本检索的关键价值。

\textbf{跨章节场景的价值。}~在 Tier-2 跨章节粗评中，
B4 MHRAG 在 Relevance、Completeness 与 Correctness 三项指标上均取得数值最高值。
其核心原因并不是 KG 替代了文本检索，
而是 KG 将 BOM 中的零件身份、CAD/IPC 中的层级关系和手册中的工序知识统一到同一实体空间，
从而缓解了单一章节关键词检索容易漏掉关系链的问题。

\textbf{Reranker 的定位。}~本文保留 BGE-Reranker-Base 作为粗排后的精排模块，
用于在 Dense、BM25 与 KG 子图之间重新选择最终上下文。
当前实验未单独量化移除 Reranker 后的三路配置表现，
因此本文不将 Reranker 的收益作为主要创新结论；
它更适合作为工程系统中的质量控制组件。

% ─────────────────────────────────────────────────────────────────────────────
\subsection{知识图谱质量对生成效果的影响}
\label{sec:disc-kg-quality}

\subsubsection{三源 KG 构建效果}

本文构建的 PT6A 整机知识图谱共 7{,}333 节点 / 8{,}752 边，
9 类核心语义关系 + 2 类辅助关系（SAME\_AS / interchangesWith）均有覆盖。
其中：
\textit{isPartOf} 2{,}869 条（BOM + CAD），
\textit{participatesIn} 1{,}750 条（手册），
\textit{precedes} 1{,}241 条（手册），
\textit{adjacentTo} 931 条（CAD 装配兄弟邻接），
\textit{specifiedBy} 649 条、\textit{matesWith} 559 条等。
本体设计涵盖完整的 7 类实体节点
（Part / Assembly / Procedure / Tool / Specification / Interface / GeoConstraint），
表明三源融合 KG 在覆盖度上相对单一文本来源的工作具有更全面的语义表达。

\subsubsection{跨源实体对齐挑战}

PT6A 维修手册原文为英文，而黄金三元组集采用中文实体名（专业术语翻译），
两者在系统检索阶段需经过 BGE-M3 多语言对齐桥接。
实验中观察到部分专业术语（如"细手锉刀"\,$\rightarrow$\,Hand File Fine）
的语言对应在 BGE-M3 默认设置下未能稳定召回，
为此本文在 QA 评测时引入双语 query 联合检索策略
（\texttt{question + question\_en} 拼接），有效缓解了此问题。
这提示工业 RAG 系统在多语种场景下应内置查询自动翻译与领域术语词典。

\subsubsection{HITL 审阅的边界价值}

第~\ref{sec:hitl}~节的 HITL 审阅机制在本文实验中并未量化测试。
其设计动机在于：
对低置信度三元组与跨源命名歧义条目的领域专家校验，
是无法被纯自动化流水线替代的。
机制本身作为 KG 构建管道的可控接口提供端到端工程价值，
其量化收益依赖于具体团队的领域专家可用度，
不在本文实验主表的核心指标范围内。

% ─────────────────────────────────────────────────────────────────────────────
\subsection{局限性与未来工作}
\label{sec:disc-limit}

\textbf{评测集规模与来源。}~Tier-1 精评集来自单一压气机子系统，
且问题由黄金三元组模板反向生成，适合验证结构化知识覆盖，
但仍不能充分代表真实生产环境中的自然提问分布。
未来应补充领域专家自由提问、多人独立标注和生产现场问答日志，
以降低模板化评测带来的偏置。

\textbf{Judge 评分的外部效度。}~Tier-2 跨章节粗评（20 题）依赖 LLM-Judge 单一指标，
不含 \texttt{gold\_chunk} 标注。
该评测方式适合验证系统覆盖广度，
但在具体题目上的判别精度受 Judge 模型自身能力影响。
未来工作可邀请领域专家对 Tier-2 答案进行二次盲审，
建立专家-Judge 一致性矩阵以校准 Judge 评分偏置。

\textbf{语料语言局限。}~本文实验语料 PT6A 维修手册为英文原文，
评测黄金集为中文专业术语翻译。
中英对应在 BGE-M3 的多语言能力下大体可工作，
但存在专业术语词典稀缺导致的部分召回失败案例。
未来工作可引入领域术语对齐词典或在检索阶段调用 LLM 进行查询改写。

\textbf{CAD 几何邻接的精度。}~由于实验环境未启用 \texttt{pythonocc-core}，
\textit{adjacentTo} 关系采用降级实现（同父装配下的兄弟节点视为逻辑邻接），
而非论文§\ref{sec:kg}~描述的轴对齐包围盒（AABB）间距计算。
该简化在装配树清晰的 PT6A IPC 文档下保持可用性，
但在复杂多层装配场景下可能产生误判。
未来工作将在生产环境集成 OCC 后回归 AABB 物理邻接计算。



% ═══════════════════════════════════════════════════════════════════════════
% 第 6 章  结论
% ═══════════════════════════════════════════════════════════════════════════
\section{结论}

% TODO

% 完成后取消注释：% ═════════════════════════════════════════════════════════════════════════════
%  第六章  结论
% ═════════════════════════════════════════════════════════════════════════════
\section{结论}
\label{sec:conclusion}

\subsection{研究总结}

针对航空发动机装配领域知识密集、来源异构、查询精确性要求高的特点，
本文提出了面向 PT6A 涡桨发动机的多源知识图谱增强混合检索生成（MHRAG）框架。
全文工作可总结为四个层面。

\textbf{在知识构建层面}，本文设计了以 BOM 解析、CAD 结构化解析与
维修手册 LLM 抽取为三轨的多智能体协同流水线，
通过三级级联实体对齐机制将三类异构数据源汇入统一的 Neo4j 知识图谱，
构建出含 7\,333 节点 / 8\,752 边、覆盖 7 类实体与 9 类关系的航空发动机装配本体。
针对 CAD 几何邻接关系，本文进一步设计了基于轴对齐包围盒（AABB）间距计算的
物理路径与基于装配树兄弟节点推断的逻辑路径，使 \textit{adjacentTo} 关系
在不同部署环境下均可获得合理覆盖。

\textbf{在检索架构层面}，本文提出三路并行检索机制，
将密集向量（BGE-M3）、稀疏检索（BM25）与
基于 BGE-M3 语义近邻搜索的知识图谱子图召回路径
通过互惠排名融合（RRF）统一排序，
并经 BGE-Reranker 交叉编码器精排送入大语言模型完成生成。
其中 Dense 负责语义召回，BM25 负责精确字符串匹配，
KG 子图负责显式提供装配层级、工序依赖和技术规范关系。

\textbf{在生成与交互层面}，本文设计了多查询扩展、
"逐步确认"模式的结构化系统提示与 SSE 流式响应的端到端生成链路，
答案附带来源片段和相关图谱关系，便于前端展示证据来源。

\textbf{在系统验证层面}，本文构建了双层 QA 评测集
（Tier-1 精评 60 题 + Tier-2 跨章节粗评 20 题）。
MHRAG 在 Tier-1 生成阶段取得 Relevance 3.50、Completeness 3.02 与
Correctness 2.35 的表现，其中 Correctness 相对 Dense-only 提升 10.8\%；
在 Tier-2 跨章节场景下，相对 BM25-only 在 Completeness 与 Correctness 维度
分别提升 8.9\% 与 1.9\%。
工程性能维度，MHRAG 在 7.76\,s 端到端延迟和 1.60\,GB 本地显存占用下
取得 35.00\% 的答非所问率，体现出面向工业现场部署的可行性。

\subsection{研究创新点}

本文相对现有工作的主要贡献可归纳为三点：

\begin{enumerate}
  \item \textbf{三源知识图谱构建}。
        相比 Liu 等\cite{liu2023avkg}的 5 类实体 / 3 类关系航空装配 KG
        与 Shi 等\cite{shi2022arkg}仅含 CAD 属性数据的 ARKG，
        本文将 BOM、CAD 与维修手册三源数据通过统一本体（7 类实体 / 9 类核心关系）
        与三级级联对齐汇入同一图谱，构建覆盖更完整的装配知识底座。

  \item \textbf{语义近邻 KG 检索}。
        相比 GraphRAG\cite{edge2024graphrag}及现有工业 RAG 工作普遍采用的
        关键词字符串匹配 KG 检索方法，
        本文将 BGE-M3 编码同时作用于查询与图谱实体描述，
        使中英混合查询能直接命中语义近邻锚点，
        降低专业术语翻译和别名表达对图谱召回的影响。

  \item \textbf{RRF 异质融合}。
        相比 HybridRAG\cite{shen2025hybridrag}的线性加权融合方案，
        本文采用互惠排名融合（RRF）在不依赖人工权重调参的前提下
        将文本块与 KG 子图统一排序，
        对单路检索失败具有更好的鲁棒性。
\end{enumerate}

\subsection{工程价值与应用前景}

本文工作的工程价值体现在三个方面。
其一，本文实现了完整的端到端工程栈
（FastAPI + Neo4j + Qdrant + React），可为航空装配领域 RAG 系统部署提供参考。
其二，分阶段 KG 构建管道（Stage 1 BOM / Stage 2 Manual / Stage 3 CAD）
配合 HITL 审阅门控，使航空领域专家可在关键节点介入校验，
契合工业场景对知识可信度的严格要求。
其三，QA 评测集（80 题）与内部基线配置可作为后续研究的对比起点。

应用前景方面，本文方法可推广至以下场景：
（1）其他型号航空发动机的装配引导系统（如 CFM56、PW1000G）；
（2）相似精度要求的工业装配领域（如重型机床、轨道交通等）；
（3）作为 MES、PLM 等制造执行系统的智能化问答前端。

\subsection{未来工作展望}

未来研究可在以下方向深化：

\textbf{评测体系扩展。}
目前 Tier-1 黄金集来自单一压气机子系统，且存在模板化问题生成带来的评测偏置。
未来将通过领域专家盲测、多人独立标注以及真实生产场景采集，
构建覆盖整机各章节的领域级 RAG 评测基准，
并尝试与公开工业问答基准建立可比性。

\textbf{KG 增量更新与版本管理。}
当前 KG 构建为离线批处理流水线，未支持手册增补、零件号变更等增量场景。
未来将引入图谱差异化更新机制，
结合 \cite{amershi2014power} 中的人在回路审阅理念，
建立面向工程现场的 KG 版本管理与回退能力。

\textbf{推理深度增强。}
当前 KG 检索路径以一跳邻域扩展为主，对涉及多步推理的复杂查询
（如"安装 X 之前需要先完成的所有工序及其工具"）尚未充分发挥图结构优势。
未来将引入基于 GraphRAG 的多跳路径搜索与
Planner-Agent 迭代探索算法\cite{shen2025hybridrag}，
进一步提升 KG 路径在复杂结构化查询上的表现。



% ═══════════════════════════════════════════════════════════════════════════
% 参考文献
% ═══════════════════════════════════════════════════════════════════════════
% 方式 A（推荐）：上传 references.bib 后启用
% \bibliographystyle{unsrtnat}
% \bibliography{references}

% 方式 B：手动列表（临时占位）
\begin{thebibliography}{99}
\bibitem{placeholder} 待补充参考文献。
\end{thebibliography}

\end{document}
